% Extracto editor-ready del contenido exclusivo de masas/08_composicion.tex.
% Las líneas 1--52 son byte-exactas a 75_rev2_ontologia_particula.tex:231--282
% y las líneas 53--396 son byte-exactas a 74_rev2_realizaciones_i.tex:641--1052
% en dos tramos contiguos. Sus propietarios íntegros permanecen preservados;
% esta residencia publica una sola vez el contenido exacto y conserva desde
% la línea 397 todo el desarrollo exclusivo del propietario de composición.

\subsection{Objet composé comme section persistante d'un produit cocyclique de parcours}
Une particule composée est une section du produit fibré des biographies constituantes, avec une frontière de liaison. Les enregistrements généalogiques sont composés avant l'application de l'extracteur de signature. Si \(\Gamma\) et \(\Lambda\) sont liés par la frontière \(B\) et l'action \(g\),

\[
(n,e)_\Gamma\star_{(g,B)}(n,e)_\Lambda
=
(n_\Gamma+n_\Lambda-n_B,
e_\Gamma+ge_\Lambda-e_B).
\]
La signature composée est

\[
\operatorname{Sig}_{\rm comp}
=L_{\rm tick}\circ\mathfrak C_{12}.
\]
Elle ne coïncide généralement pas avec la somme des signatures individuelles, car l'enregistrement généalogique de liaison porte énergie, retenue et mémoire. C'est là que réside l'énergie de liaison.

\Needspace{7\baselineskip}
\subsection{Résonance comme section transitoire et largeur comme taux de désintégration}
Une résonance est une section métastable dont la résolvante possède un pôle d'énergie complexe. Si \(z_X=M_Xc_{\rm int}^{2}-i\Gamma_X/2\), la masse est \(M_X=c_{\rm int}^{-2}\operatorname{Re}z_X\) ; la largeur mesure sa durée de vie moyenne :

\[
\Gamma_X=\frac{\hbar_{\rm int}}{\tau_X^{\rm life}}.
\]
La largeur n'est pas une sixième coordonnée de \(\sigma\). Elle procède de l'opérateur de survie, des canaux de désintégration et du prolongement analytique de la résolvante. L'inventaire de 384 largeurs exige donc un lecteur propre, distinct de celui de l'inventaire de 471 masses.

\Needspace{7\baselineskip}
\subsection{Section virtuelle comme extension intermédiaire non asymptotique de l'atlas}
Une section virtuelle est une extension finie qui participe à une composition ou à un propagateur sans définir de branche globale stable dans le codomaine physique. Si \(L(s)\) est son horizon,

\[
L(s)<\infty,
\qquad
\operatorname{Ext}_{L(s)}(s)\ne\varnothing,
\qquad
\operatorname{Ext}_{L(s)+1}(s)=\varnothing.
\]
La virtualité est une obstruction terminale, et non une étiquette documentaire ou une masse imaginaire. Une résonance peut être instable sans être virtuelle en ce sens ; une ligne interne peut être virtuelle sans constituer une nouvelle espèce.

\Needspace{7\baselineskip}
\subsection{Section de Majorana comme point fixe de la conjugaison particule–antiparticule}
Quatre constructions doivent rester distinctes. L'involution des mots \(C_M\) définit l'antisection ; un spineur de Majorana satisfait une condition de réalité ; un mode zéro de Majorana est une excitation localisée d'énergie nulle dans un Hamiltonien topologique ; le secteur d'Ising possède les règles de fusion

\[
\sigma\otimes\sigma=1\oplus\psi.
\]
Une particule physique est de Majorana lorsque sa représentation fermionique admet une autoconjugaison compatible avec le Hamiltonien, la localisation et le produit scalaire. La neutralité électrique et l'invariance d'un mot sous \(C_M\) ne suffisent pas. En particulier, déterminer la nature Majorana des neutrinos exige un entrelaceur entre l'opérateur neutre, la base des masses et la structure réelle fermionique.

\Needspace{7\baselineskip}
\subsection{Section anyonique et représentation non triviale du groupe de tresses}
Un anyon est une section bidimensionnelle dont la représentation du groupe de tresses ne se réduit pas à la permutation bosonique ou fermionique. Dans le secteur de Fibonacci,

\[
\tau\otimes\tau=1\oplus\tau,
\qquad
\operatorname{FPdim}(\tau)=\varphi.
\]
Dans le secteur abélien \(\mathbb Z/6\mathbb Z\), le quotient d'enroulement possède six caractères unidimensionnels. Pour obtenir une particule physique, il faut des matrices \(F\) et \(R\), un pentagone, un hexagone, un Hamiltonien à gap et des opérateurs de tressage. L'anneau de fusion exact détermine le codomaine topologique, mais ne fixe pas, à lui seul, une masse universelle.

\Needspace{7\baselineskip}
\subsection{Classes tachyonique et fantôme comme obstructions de positivité et de stabilité}
Une particule persistante possède une évolution unitaire, une énergie bornée inférieurement et une masse spectrale réelle sur la droite physique. Un tachyon exigerait un pôle avec \(m^2<0\) demeurant un état asymptotique stable. Dans l'architecture HMT--MD, une direction qui quitte le cône des réalisations positives revient à la frontière nulle, perd sa survie ou se manifeste comme l'instabilité d'un fond qui doit se condenser. Elle ne définit pas de section persistante supraluminique. L'impossibilité se formule comme une incompatibilité entre persistance globale, positivité spectrale et \(m^2<0\), et non comme une interdiction des paramètres tachyoniques transitoires dans un développement perturbatif.

Un fantôme est un mode de norme négative. La positivité du produit scalaire et la conservation unitaire excluent sa réalisation physique. Les fantômes de Faddeev--Popov appartiennent au complexe de jauge et non à l'espace physique des états ; cette distinction doit être conservée.

\Needspace{7\baselineskip}
\section{Composition, résonances et largeurs}
\Needspace{7\baselineskip}
\subsection{Composition non additive des masses par la loi cocyclique de liaison}
La Loi Générale des Masses agit sur des biographies généalogiques. Lorsque plusieurs constituants forment un état lié, l'objet d'entrée n'est pas la liste de leurs masses isolées, mais un enregistrement composite qui conserve l'ordre du couplage, la frontière commune, la retenue, l'orientation et la mémoire. Soient

\[
\Gamma_1,\ldots,\Gamma_n
\]
des parcours constituants, \(\mathcal T\) un arbre temporel de composition et \(\partial\) la frontière de liaison. L'opérateur dodécaphasique de composition s'écrit

\[
\mathfrak C_{12}^{(n)}:
(\Gamma_1,\ldots,\Gamma_n;\mathcal T,\partial)
\longmapsto \mathcal L_{\rm comp}.
\]
Ce n'est qu'ensuite que s'appliquent le lecteur de l'enregistrement, l'extracteur de signature, le caractère et les tours :

\[
\mathcal L_{\rm comp}
\longmapsto \sigma_{\rm comp}
\longmapsto \chi_{\rm comp}
\longmapsto \widehat M_{\rm comp}.
\]
Cette priorité est essentielle. Si les masses isolées étaient d'abord évaluées puis additionnées, les données qui distinguent un état lié d'un seuil de constituants libres disparaîtraient précisément. L'énergie de liaison n'est pas un terme ajouté à la fin : elle est l'image dimensionnelle d'une différence d'enregistrement produite par la liaison.

\Needspace{7\baselineskip}
\subsection{Loi cocyclique de liaison}
Pour deux enregistrements \(L_1,L_2\), la composition prend la forme

\[
L_1\star L_2
=L_1+L_2+\mathfrak b(L_1,L_2),
\]
où \(\mathfrak b\) enregistre frontière, retenue et mémoire partagées. L'indépendance à l'égard du parenthésage exige l'identité cocyclique

\[
\mathfrak b(L_1,L_2)
+\mathfrak b(L_1\star L_2,L_3)
=
\mathfrak b(L_2,L_3)
+\mathfrak b(L_1,L_2\star L_3).
\]
Si cette égalité échoue, les deux façons de construire un même arbre triple produisent des enregistrements généalogiques distincts et il n'existe pas d'espèce composée bien définie. Si elle est satisfaite, le cocycle ne disparaît pas : il garantit que sa contribution est une propriété de l'état et non un artefact de notation.

Une présentation plus explicite conserve en outre l'action de bord. Soit \(e\in\mathbb Z^{12}\) le vecteur ordonné des événements, et \(n\) le compteur d'action ou d'événements d'un enregistrement. Si \((n,e)_\Gamma\) et \((n,e)_\Lambda\) sont des enregistrements, si \(B\) est le sous-enregistrement commun de données \((n_B,e_B)\), et si \(g\in\operatorname{Aut}(\mathbb Z^{12})\) transporte l'orientation du second enregistrement vers celle du premier en préservant le lecteur,

\[
(n,e)_\Gamma\star_{(g,B)}(n,e)_\Lambda
=
(n_\Gamma+n_\Lambda-n_B,
e_\Gamma+g e_\Lambda-e_B).
\]
Le terme \(B\) évite de compter deux fois la frontière identifiée ; l'action \(g\) transporte le second enregistrement dans l'orientation du premier. La signature composée est donc

\[
\operatorname{Sig}_{\rm comp}
=L_{\rm tick}\!\left(
\mathfrak C_{12}^{(n)}
(\Gamma_1,\ldots,\Gamma_n;\mathcal T,\partial)
\right),
\]
et non la somme automatique des signatures élémentaires.

\Needspace{7\baselineskip}
\subsection{Clôture cocyclique de la composition mésonique quark–antiquark}
Une expression mésonique orientée prend la forme

\[
M(q,q')=q\otimes(q')^\vee.
\]
Le domaine brut contient \(6\times6\times3\times4=432\) expressions lorsque l'on considère saveur, couleur compatible et orientation. Cette cardinalité n'est pas le nombre de mésons physiques. Avant d'individualiser une espèce, on applique :

\begin{enumerate}
\item le projecteur de singulet de couleur ;
\item le projecteur de charge et d'isospin ;
\item la représentation de spin et de parité ;
\item l'arbre temporel de liaison ;
\item la condition radiale ou orbitale ;
\item le canal de préparation et de désintégration lorsque l'état est résonant.
\end{enumerate}
Le descripteur physique complet est

\[
\mathcal I_M=
(q\bar q',J^{PC},I,I_3,S,C,B,\mathcal T,\partial,
\Gamma_M,\mathcal L_M,\sigma_M,\eta_M).
\]
Deux états de mêmes saveurs peuvent différer par \(J^{PC}\), l'excitation radiale, le moment orbital, le feuillet ou le canal, et posséder ainsi des pôles distincts. La masse ne se déduit pas du seul nom des constituants.

Les évaluations publiées pour \(\pi^0,\pi^\pm,K^\pm,K^0\) sont exactes une fois fixés les enregistrements historiques déclarés. L'affirmation plus forte --- la génération prospective de chaque enregistrement nominal à partir de son arbre de liaison --- exige que \(\mathfrak C_{12}\) reconstruise ces enregistrements sans consulter la valeur qui sera comparée. Les deux couches sont conservées ensemble : le calcul conditionnel n'est pas supprimé, mais n'est pas non plus présenté comme un sélecteur d'identité.

\Needspace{7\baselineskip}
\subsection{Clôture trilinéaire de la composition baryonique et neutralisation du caractère de couleur}
Une expression baryonique orientée utilise trois constituants et trois couleurs. Le domaine brut du dénombrement adopté contient

\[
\lvert\mathcal B_{\rm RGB}^{\rm raw}\rvert=1728
\]
expressions. La projection physique doit imposer simultanément :

\[
P_{\rm singlete\ de\ color},
\qquad
P_{J^P},
\qquad
P_{I,I_3,S,C,B},
\qquad
P_{\rm Pauli}.
\]
L'arbre triple conserve la paire qui s'est couplée en premier et la façon dont le troisième constituant a traversé la frontière. L'identité cocyclique garantit que les parenthésages équivalents représentent le même état lorsque l'entrelaceur physique les identifie ; les canaux non équivalents demeurent des états distincts.

Les signatures historiques

\[
\sigma_p=(59,0,9,1,0),
\qquad
\sigma_n=(59,1,-34,0,1)
\]
produisent les évaluations publiées du proton et du neutron lorsque l'on applique le caractère complet. Ce sont des contrôles exacts du lecteur pour ces enregistrements. L'explication physique de la différence neutron--proton exige en outre de dériver, avant la confrontation, le signe de l'isospin, l'orientation des lacunes, le bord de liaison et la contribution électromagnétique dans une carte commune.

\Needspace{7\baselineskip}
\subsection{Réalisations composées exotiques et conditions de stabilité de la liaison}
La même grammaire admet davantage de constituants sans modifier la loi :

\[
\mathfrak C_{12}^{(n)}:
\prod_{j=1}^{n}\Gamma_j\times\mathcal T_n\times\partial_n
\longrightarrow\mathcal L_{\rm comp}.
\]
Un tétraquark correspond à \(qq\bar q\bar q\) ; un pentaquark, à \(qqqq\bar q\) ; une boule de glu, à des parcours de la représentation adjointe de couleur projetés sur le singulet ; un hybride, à des parcours fermioniques et de jauge réunis dans un même arbre. La liste des constituants ne détermine pas, à elle seule, l'espèce. Une classification physique complète doit inclure couleur, spin, parité, conjugaison, isospin, topologie de l'arbre et pôle de l'opérateur couplé.

Cette formulation distingue aussi deux descriptions concurrentes d'un même candidat --- molécule hadronique et état compact, par exemple --- : ce sont deux arbres ou deux projecteurs sur un même secteur de constituants. Elles ne deviennent équivalentes que s'il existe un opérateur unitaire entrelaçant leurs opérateurs, leurs canaux et leurs pôles.

\Needspace{7\baselineskip}
\subsection{État stable, état lié et résonance}
Il convient de distinguer trois notions.

Un état stable est un vecteur propre normalisable du Hamiltonien physique, d'énergie réelle et sans canal de désintégration ouvert. Un état lié est un vecteur propre situé sous le seuil pertinent ; il peut être stable ou se désintégrer par un autre canal. Une résonance est un pôle du prolongement de la résolvante ou de la matrice de diffusion. Elle n'est pas un vecteur propre normalisable du Hamiltonien autoadjoint sur l'axe réel.

Pour un pôle d'énergie

\[
z_X=E_X-\frac{i}{2}\Gamma_X
=M_Xc_{\rm int}^{2}-\frac{i}{2}\Gamma_X,
\qquad \Gamma_X>0,
\]
l'énergie du pôle est \(E_X=\operatorname{Re}z_X\), sa masse est \(M_X=c_{\rm int}^{-2}\operatorname{Re}z_X\) et sa largeur énergétique est \(\Gamma_X=-2\operatorname{Im}z_X\). La Loi Générale des Masses fournit la branche réelle conditionnée par l'identité et l'enregistrement ; le prolongement ouvert des canaux détermine le déplacement complexe et la largeur. Dans une carte d'unités naturelles \(c_{\rm int}=1\), les coordonnées numériques d'énergie et de masse coïncident, mais leurs types ne sont pas identifiés avant le choix de cette carte.

\Needspace{7\baselineskip}
\subsection{Correction du signe dans la lecture discrète du pôle}
Soit \(t_0>0\) un pas temporel et supposons que le mode d'un pas ait pour valeur propre

\[
\lambda=r e^{-i\theta},
\qquad 0<r<1,
\]
avec la convention

\[
\lambda=\exp\!\left(-\frac{i}{\hbar_{\rm int}}z t_0\right).
\]
Alors

\[
z=\frac{\hbar_{\rm int}}{t_0}
\bigl(\theta+i\log r\bigr)
=E-\frac{i}{2}\Gamma
=Mc_{\rm int}^{2}-\frac{i}{2}\Gamma,
\]
de sorte que

\[
E=\frac{\hbar_{\rm int}}{t_0}\theta,
\qquad
M=\frac{\hbar_{\rm int}}{t_0c_{\rm int}^{2}}\theta,
\qquad
\Gamma=-\frac{2\hbar_{\rm int}}{t_0}\log r>0.
\]
Le signe positif devant \(i\log r\) est imposé : puisque \(\log r<0\), la partie imaginaire de \(z\) est négative. La forme \(\theta-i\log r\) placerait le pôle dans le demi-plan supérieur et décrirait une croissance, et non une décroissance.

La probabilité de survie après \(N\) pas est

\[
|\lambda|^{2N}=r^{2N}
=\exp\!\left(-\frac{\Gamma}{\hbar_{\rm int}}Nt_0\right).
\]
La durée de vie associée est donc

\[
\tau_X^{\rm life}=\frac{\hbar_{\rm int}}{\Gamma_X}.
\]
Cette identité démontre aussi pourquoi la largeur ne doit pas être introduite comme sixième coordonnée de la signature \(\mathbb Z^5\) : elle naît du module du mode ouvert, et non du caractère positif qui organise la masse.

\Needspace{7\baselineskip}
\subsection{Opérateur effectif des canaux}
Soit \(P\) le sous-espace préparé et \(Q=I-P\) l'ensemble des canaux éliminés. Le complément est incorporé au moyen de l'opérateur de Feshbach

\[
H_{\rm eff}(z)
=PHP+PHQ(z-QHQ)^{-1}QHP.
\]
Les pôles satisfont

\[
\det\bigl(z-H_{\rm eff}(z)\bigr)=0.
\]
Le premier terme conserve l'opérateur interne préparé. Le second est l'autoénergie des canaux : il déplace la partie réelle et produit une partie imaginaire lorsque le continu s'ouvre. La masse du caractère et la largeur de survie convergent vers le même pôle, mais procèdent d'applications distinctes.

Une formulation équivalente utilise un projecteur de survie distinct du complément de Feshbach. Soit

\[
S_X=S_X^*=S_X^2
\]
le projecteur sur le sous-espace survivant de l'état préparé. L'opérateur comprimé d'un pas est

\[
T_X=S_XUS_X.
\]
Supposons que \(T_X\) possède une valeur propre dominante simple \(\lambda_X\), que l'état préparé ait un recouvrement non nul avec son sous-espace spectral et que le reste du spectre soit séparé en module. Si \(r_X=|\lambda_X|=\rho(T_X)<1\), alors

\[
\Gamma_X=-\frac{2\hbar_{\rm int}}{t_0}\log r_X.
\]
Si l'opérateur est normalisé dès le départ sur les probabilités plutôt que sur les amplitudes, le facteur deux disparaît. Le texte doit déclarer laquelle des deux conventions il utilise ; les mélanger produit une largeur erronée d'un facteur deux.

\Needspace{7\baselineskip}
\subsection{Conjugaison de la particule et de l'antiparticule}
Soit \(C_M\) l'involution qui inverse orientation et feuillet, et \(U_C\) sa réalisation unitaire sur l'espace physique. L'égalité des masses de pôle et des largeurs exige la covariance de l'opérateur comme des canaux :

\[
U_CH_{\rm eff}^{X}(z)U_C^{-1}
=H_{\rm eff}^{\bar X}(z).
\]
Alors les ensembles de pôles retardés coïncident après transport du secteur et de la prescription causale, et

\[
M_{\bar X}=M_X,
\qquad
\Gamma_{\bar X}=\Gamma_X.
\]
La seule équivalence unitaire ne transforme pas le pôle retardé en son conjugué complexe ; elle transporte le même problème spectral vers le secteur antiparticule. La parité de la signature massique ne suffit pas à démontrer la seconde égalité : les canaux et leurs mesures doivent eux aussi être covariants. Si la symétrie est représentée antiunitairement, il faut également transporter la prescription retardée ; la seule conjugaison \(z\mapsto\bar z\) échange pôles retardés et avancés et ne constitue pas, à elle seule, l'affirmation physique de l'égalité des durées de vie moyennes.

\Needspace{7\baselineskip}
\subsection{Inventaire typé des réalisations et des observables de comparaison}
La coupe externe conservée contient

\[
471\ \text{observables de masse},
\qquad
384\ \text{observables de largeur}.
\]
La partition des masses est

\Needspace{7\baselineskip}
\begingroup
\footnotesize
\setlength{\tabcolsep}{1.9pt}
\begin{longtable}{L{5.04cm}L{5.04cm}L{5.04cm}}
\toprule
\rowcolor{HMTNavy}\HMTtablehead{Domaine physique d'arrivée} & \HMTtablehead{Masses} & \HMTtablehead{Largeurs} \\
\midrule
\endfirsthead
\multicolumn{3}{l}{\sffamily\scriptsize\color{HMTGray}Suite du tableau}\\[-1pt]
\toprule
\rowcolor{HMTNavy}\HMTtablehead{Dominio físico de llegada} & \HMTtablehead{Masas} & \HMTtablehead{Anchuras} \\
\midrule
\endhead
\midrule
\endfoot
\bottomrule
\endlastfoot
partículas elementales, leptónicas y extensiones & 53 & 9 \\
\rowcolor{HMTLight}quarks y extensiones quark & 8 & 1 \\
mesones & 243 & 225 \\
\rowcolor{HMTLight}bariones & 166 & 149 \\
registro gravitatorio de contraste & 1 & 0 \\
\rowcolor{HMTLight}\textbf{Total} & \textbf{471} & \textbf{384} \\
\end{longtable}
\endgroup
Estas filas son observables, no especies elementales. Varias filas pueden referirse a diferentes cargas de un mismo multiplete, estimadores de una misma resonancia, masas Breit--Wigner, masas de polo o estados excitados. La individualización interna precede siempre a la fila externa.

El inventario conserva además una propiedad metodológica fuerte: ninguna de sus 855 magnitudes fue autorizada como selector de ruta. Los 471 registros de masa y los 384 de anchura sólo se abren después de congelar la identidad HMT, el proyector, el operador y la carta de comparación.

\iffalse
\Needspace{7\baselineskip}
\subsection{Malla histórica extendida de pares y resonancias}
Una formulación histórica asoció a la malla angular extendida la expresión

\[
\Theta(n,k,t)
=nA_{\rm rad}-k\frac{C^*_{\rm rad}}6+t\Delta_4.
\]
Avec les définitions actuelles de \(A\), \(C^*\), \(\Delta_4\) et de la carte exponentielle, cette expression isolée ne reproduit pas les valeurs archivées. Elle est conservée comme hypothèse historique de paramétrisation, et non comme leur générateur exécutable. Nous conservons les douze évaluations importées parce qu'elles font partie de la généalogie numérique de la loi. Cinq noms --- \(\eta,\rho^0,\phi(1020),J/\psi\) et \(\Upsilon(1S)\) --- élargissent concrètement les étiquettes visibles par rapport aux dix-sept évaluations de référence :

\Needspace{7\baselineskip}
\begingroup
\fontsize{7.2}{8.2}\selectfont
\setlength{\tabcolsep}{1.9pt}
\begin{longtable}{L{1.25cm}L{0.78cm}L{0.78cm}L{0.78cm}L{3.43cm}L{3.90cm}L{4.06cm}}
\toprule
\rowcolor{HMTNavy}\HMTtablehead{État} & \HMTtablehead{\(n\)} & \HMTtablehead{\(k\)} & \HMTtablehead{\(t\)} & \HMTtablehead{Référence ultérieure (MeV/\allowbreak{}c\({}^{2}\))} & \HMTtablehead{Évaluation HMT historique (MeV/\allowbreak{}c\({}^{2}\))} & \HMTtablehead{Diferencia (ppm)} \\
\midrule
\endfirsthead
\multicolumn{7}{l}{\sffamily\scriptsize\color{HMTGray}Suite du tableau}\\[-1pt]
\toprule
\rowcolor{HMTNavy}\HMTtablehead{Estado} & \HMTtablehead{\(n\)} & \HMTtablehead{\(k\)} & \HMTtablehead{\(t\)} & \HMTtablehead{Referencia posterior (MeV/\allowbreak{}c\({}^{2}\))} & \HMTtablehead{Evaluación HMT histórica (MeV/\allowbreak{}c\({}^{2}\))} & \HMTtablehead{Diferencia (ppm)} \\
\midrule
\endhead
\midrule
\endfoot
\bottomrule
\endlastfoot
\(\mu\) & 64 & 146 & 1 & 105.658374 & 105.656529 & -17.461938\allowbreak{}2274 \\
\rowcolor{HMTLight}\(\tau\) & 64 & 349 & -1 & 1776.86 & 1776.835 & -14.069763\allowbreak{}5154 \\
\(\pi^\pm\) & 37 & 125 & 0 & 139.57039 & 139.570582 & 1.375649\allowbreak{}94982 \\
\rowcolor{HMTLight}\(K^\pm\) & 52 & 177 & 1 & 493.677 & 493.679 & 4.051231\allowbreak{}87833 \\
\(K^0\) & 53 & 178 & 1 & 497.611 & 497.616 & 10.048009\allowbreak{}3889 \\
\rowcolor{HMTLight}\(p\) & 55 & 195 & 1 & 938.272081 & 938.27244 & 0.382618\allowbreak{}226919 \\
\(n\) & 55 & 197 & 0 & 939.565413 & 939.538 & -29.176254\allowbreak{}9161 \\
\rowcolor{HMTLight}\(\eta\) & 46 & 158 & 0 & 547.862 & 547.836 & -47.457206\allowbreak{}3768 \\
\(\rho^0\) & 50 & 176 & 0 & 775.26 & 775.267 & 9.029228\allowbreak{}90385 \\
\rowcolor{HMTLight}\(\phi(1020)\) & 54 & 190 & 0 & 1019.461 & 1019.475 & 13.732747\allowbreak{}0104 \\
\(J/\allowbreak{}\psi\) & 64 & 237 & 0 & 3096.916 & 3096.979 & 20.342818\allowbreak{}4684 \\
\rowcolor{HMTLight}\(\Upsilon(1S)\) & 89 & 334 & 0 & 9460.3 & 9460.34 & 4.228195\allowbreak{}72318 \\
\end{longtable}
\endgroup
Las doce cifras archivadas son evaluaciones históricas importadas; no se atribuyen a \(\Theta\) hasta localizar la carta que las reproduzca. Además, la malla no construye el morfismo

\[
\text{état composé}
\longrightarrow\text{route et registre généalogique}
\longrightarrow(n,k,t).
\]
Elle conserve en outre la référence externe dans le même tableau historique. Sa force probante est donc celle d'une preuve numérique conditionnelle, et non d'une sélection physique antérieure à la confrontation. La présenter élargit le dénombrement visible sans transformer cinq résonances en cinq prédictions antérieures à la confrontation, ni supplanter les évaluations E3 en vigueur.
\fi

\Needspace{7\baselineskip}
\subsection{Normalisation des identités nominales en enregistrements comparables}
Le raccordement prospectif complet prend la forme

\[
\begin{aligned}
\text{identité physique}
&\longrightarrow
\text{secteur et représentation}
\longrightarrow
\text{projecteur de fibre ou d'arbre composé}\\
&\longrightarrow
\text{route enrichie}
\longrightarrow
\text{registre généalogique}
\longrightarrow
\text{signature et tours}\\
&\longrightarrow
\text{opérateur préparé}
\longrightarrow
\text{pôle ou valeur propre}
\longrightarrow
\text{carte dimensionnelle}\\
&\longrightarrow
\text{comparaison externe}.
\end{aligned}
\]
Pour un état stable, le résultat opératoriel est une valeur propre réelle. Pour une résonance, c'est un pôle complexe. Pour un quark, l'application inclut schéma et échelle. Pour une limite expérimentale, elle inclut le niveau de confiance. Pour un observable corrélé, elle inclut la matrice de covariance. Ces différences sont structurelles : elles déterminent quelles grandeurs peuvent être comparées.

\Needspace{7\baselineskip}
\subsection{Conditions de clôture individuelle des réalisations matérielles}
Une ligne de l'inventaire est individualisée à partir de HMT--MD lorsque sont construits, sans consulter sa grandeur externe :

\begin{enumerate}
\item le secteur et la représentation ;
\item le parcours, la multisection ou l'arbre de composition ;
\item l'enregistrement généalogique et la signature ;
\item l'opérateur de masse et, le cas échéant, l'opérateur des canaux ;
\item le projecteur qui sélectionne l'état nommé ;
\item la convention de masse, de largeur, de schéma et d'échelle ;
\item la carte dimensionnelle ;
\item le critère interne de réfutation.
\end{enumerate}
La loi universelle fournit déjà la grammaire commune. Le travail par état ne consiste pas à inventer une formule différente, mais à construire l'entrelaceur qui transporte une identité physique jusqu'à l'opérateur approprié.

\Needspace{7\baselineskip}
\subsection{obstructions et conditions spécifiques de non-dégénérescence}
La composition est réfutée si deux parenthésages physiquement équivalents donnent des signatures ou des spectres distincts. Le sélecteur nominal est réfuté s'il change lorsque la colonne externe de masse est masquée. Le lecteur de largeur est réfuté s'il produit \(\Gamma<0\) pour une contraction, confond amplitude et probabilité ou attribue une largeur à une section stable sans canal ouvert. La carte de résonance est réfutée si elle ne déclare pas masse de pôle ou de Breit--Wigner. La comparaison des quarks reste indéterminée si elle omet le schéma ou l'échelle. L'égalité des largeurs particule--antiparticule est réfutée si l'opérateur des canaux n'est pas covariant par conjugaison.

\Needspace{7\baselineskip}
\subsection{Transport des classes matérielles vers le corridor exceptionnel au moyen des parcours, signatures et représentations}
La grammaire de composition, les domaines mésonique et baryonique, l'opérateur de Feshbach, le dénombrement 471/384 et le principe de comparaison sont conservés dans le développement intégral antérieur, sections 16 et 18. Les enregistrements individualisés sont présentés intégralement dans les inventaires exhaustifs de masses et de largeurs de l'appendice. La correction du signe du pôle découle directement de \(\lambda=\exp(-izt_0/\hbar_{\rm int})\) et renforce la formulation sans modifier aucun de ces enregistrements.

\Needspace{7\baselineskip}
\section{Secteurs topologiques, extensions hypothétiques et gravitation}
\Needspace{7\baselineskip}
\subsection{Séparation typée entre secteurs topologiques, extensions hypothétiques et réalisation gravitationnelle}
L'architecture distingue précisément :

\begin{enumerate}
\item une identité algébrique interne ;
\item un codomaine de fusion ou de tressage ;
\item une particule physique avec Hamiltonien, gap, localisation et canal de
mesure.

\end{enumerate}
Une règle de fusion exacte n'est pas encore une masse. Un mode de bord n'est pas encore une espèce. Une involution de parcours n'est pas encore une condition de Majorana. Le passage physique exige un foncteur qui préserve produit, involution, orientation, énergie et composition des chemins.

\Needspace{7\baselineskip}
\subsection{Secteur de Fibonacci : anneau basé, fusion et représentation de tresses}
L'incidence surfacique--volumétrique est régie par

\[
F_{\rm av}=
\begin{pmatrix}
0&1\\[2pt]1&1
\end{pmatrix}.
\]
Son polynôme caractéristique est

\[
x^2-x-1,
\]
et sa valeur propre dominante est \(\varphi\). Dans la base \(\{\mathbf1,\tau\}\), la multiplication par \(\tau\) est représentée par \(F_{\rm av}\) ; ainsi

\[
\tau\otimes\tau\cong\mathbf1\oplus\tau,
\qquad
\operatorname{FPdim}(\tau)=\varphi.
\]
C'est une réalisation exacte de l'anneau basé de Fibonacci. Pour obtenir une catégorie anyonique physique, il faut en outre construire l'associateur \(F\), le tressage \(R\), les équations du pentagone et de l'hexagone, la chiralité, un Hamiltonien local à gap et l'application de préparation et de lecture. Le résultat de masse, lorsqu'il existe, est un gap d'énergie converti au moyen d'une vitesse effective :

\[
m_{\rm brecha}=\frac{\Delta E}{c_{\rm eff}^{2}}.
\]
Il n'existe pas de masse universelle de l'anyon de Fibonacci déductible uniquement de \(\varphi\).

\Needspace{7\baselineskip}
\subsection{Séparation entre le secteur topologique et le centre pointé de l'atlas}
Le centre tressé du secteur pointé fondé sur \(\mathbb Z/9\mathbb Z\) possède 81 objets simples inversibles, chacun de dimension quantique un. Un objet simple de Fibonacci exigerait la dimension \(\varphi\), et un objet simple d'Ising exigerait \(\sqrt2\). Aucun des deux n'apparaît donc comme objet simple interne de ce centre pointé.

Ce contrôle n'élimine pas les réalisations topologiques. Il détermine leur place : elles doivent s'obtenir par extension, condensation, quotient, défaut ou foncteur monoïdal vers un codomaine non pointé. Présenter l'anneau de Fibonacci comme s'il constituait déjà toute la catégorie effacerait précisément le morphisme physique à construire.

\Needspace{7\baselineskip}
\subsection{Secteur d'Ising et modes de Majorana}
La catégorie d'Ising a pour objets simples

\[
\mathbf1,\quad\psi,\quad\sigma
\]
et règles

\[
\psi\otimes\psi=\mathbf1,
\qquad
\psi\otimes\sigma=\sigma,
\qquad
\sigma\otimes\sigma=\mathbf1\oplus\psi,
\qquad
d_\sigma=\sqrt2.
\]
Quatre objets distincts reçoivent souvent des noms proches et doivent demeurer séparés :

\begin{itemize}
\item l'involution des parcours \(C_M=-\operatorname{rev}\) ;
\item la parité de la complétion CAR ;
\item un spineur relativiste autoconjugué ;
\item un mode zéro topologique de Majorana dans une phase d'Ising.
\end{itemize}
Un fermion relativiste est de Majorana lorsqu'il existe une conjugaison réelle \(\mathcal C\) compatible avec sa représentation et sa dynamique,

\[
\psi=\mathcal C\bar\psi^{T},
\]
et que le terme de masse respecte cette condition. Un mode zéro de Majorana se décrit par un opérateur autoadjoint \(\gamma=\gamma^\dagger\) qui satisfait

\[
\{\gamma_i,\gamma_j\}=2\delta_{ij}
\]
et reste séparé du continu par un gap. L'égalité d'un mot avec son image par \(C_M\) ne démontre aucune de ces deux réalisations.

Pour les neutrinos, l'alternative Dirac/Majorana ne se décide qu'après construction de l'opérateur neutre, de sa base de masses, de l'entrelaceur de mélange et de la structure réelle fermionique. La neutralité électrique est nécessaire à une autoconjugaison non triviale, mais ne suffit pas.

\Needspace{7\baselineskip}
\subsection{Tressage abélien d'ordre \(6\)}
Soit

\[
\operatorname{wind}_6:B_n\longrightarrow\mathbb Z/6\mathbb Z
\]
le lecteur d'enroulement et

\[
\chi_k(b)=\exp\!\left(
\frac{2\pi i k}{6}\operatorname{wind}_6(b)\right).
\]
Les réalisations de phase \(q\) se réduisent à de simples permutations seulement lorsque

\[
q^2=1.
\]
Les quatre caractères restants distinguent des classes algébriques du lecteur d'enroulement. Le résultat exact est la classification des six caractères unidimensionnels du quotient d'enroulement ; ce n'est pas encore une classification de six phases matérielles. Affirmer qu'ils conservent un tressage physique exige des opérateurs de tresse et une représentation du groupe de tresses. Leur existence n'implique pas quatre nouvelles particules : une espèce anyonique exige un Hamiltonien, des opérateurs de tresse, un gap, une stabilité et un couplage à une multisection physique. Leur promotion au statut de phases matérielles dépend de cette réalisation.

\Needspace{7\baselineskip}
\subsection{Positivité et exclusion du tachyon persistant}
La droite d'action possède une orientation positive,

\[
\hbar_{\rm int}>0,
\]
le caractère de masse prend des valeurs positives,

\[
\chi_{\rm full}(\sigma,\eta)>0,
\]
et le raffinement bidirectionnel est réalisé par conjugaison. Supposons explicitement

\[
\widehat M^{(0)}\ge0,
\qquad
\widehat K=\widehat K^*,
\qquad
R_{\rm act}>0.
\]
Alors

\[
\widehat M^{\beta}
=R_{\rm act}^{\widehat K/2}
\widehat M^{(0)}
R_{\rm act}^{\widehat K/2}\ge0.
\]
En conséquence,

\[
\operatorname{Spec}(\widehat M^\beta)\subset[0,\infty).
\]
Ainsi, l'opérateur de masse au carré

\[
\widehat{M^2}_{\beta}:=(\widehat M^\beta)^2
\]
est positif et satisfait

\[
\operatorname{Spec}(\widehat{M^2}_{\beta})\subset[0,\infty).
\]
Une particule asymptotique persistante satisfaisant les prémisses de positivité et d'unitarité ne peut avoir \(m^2<0\). Si la linéarisation d'un fond produit une fréquence imaginaire, l'objet est une direction instable : le fond doit se relaxer ou se condenser. Si le parcours cesse de se prolonger, l'objet est une section virtuelle à obstruction terminale. Aucun des deux cas ne définit une particule stable supraluminique.

La portée de la proposition est fixée par ses prémisses. Un paramètre « tachyonique » dans un potentiel perturbatif peut signaler une transition de phase ; ce qui est exclu est sa promotion au statut d'espèce persistante de norme positive.

\Needspace{7\baselineskip}
\subsection{Exclusion indépendante des fantômes}
Un fantôme physique serait un mode de norme négative. Soit \(\mathcal H_{\rm phys}=\ker G/\operatorname{im}G_0\) le quotient de jauge muni d'un produit scalaire positif, et \(\widehat M_{\rm phys}\) autoadjoint sur ce quotient. Alors la mesure spectrale d'un état préparé satisfait

\[
\mu_{P,a}^{\rm mass}(B)\ge0.
\]
Une direction négative indique que la positivité du quotient a échoué ou que le mode n'appartient pas à l'espace physique ; elle ne disparaît pas par une affirmation nominale. Les fantômes de Faddeev--Popov sont des variables du complexe de jauge, et non des particules asymptotiques. Cette exclusion est distincte de l'exclusion tachyonique : un fantôme viole la positivité de la norme ; une direction tachyonique viole la stabilité du fond.

\Needspace{7\baselineskip}
\subsection{Photon, biais unidirectionnel et lien Proca--Stückelberg}
Les canaux harmoniques produisent le biais

\[
\varepsilon=q_-^{90}-q_+^{90}.
\]
Une réalisation cinématique bidirectionnelle peut s'écrire

\[
c_+=\frac{c_{\rm int}}{1+\varepsilon},
\qquad
c_-=\frac{c_{\rm int}}{1-\varepsilon},
\]
pour laquelle la mesure d'aller et retour conserve exactement

\[
\frac{2}{c_+^{-1}+c_-^{-1}}=c_{\rm int}.
\]
Le biais demeure dans la lecture unidirectionnelle et s'annule dans la lecture réciproque. Cette identité appartient à l'architecture à deux orientations ; elle ne démontre pas, à elle seule, une masse photonique.

Si la clôture de frontière produit une fréquence de gap \(\mu_T\in\mathcal L_T^{-1}\) dans l'enregistrement,

\[
\omega^2=c_{\rm int}^{2}k^2+\mu_T^2,
\]
un entrelaceur dimensionnel peut la transporter vers une réalisation Proca--Stückelberg de masse effective

\[
m_\gamma=\frac{\hbar_{\rm int}\mu_T}{c_{\rm int}^{2}}.
\]
Pour une propagation de fréquence \(\omega\), la correction temporelle du premier ordre satisfait

\[
\frac{\Delta T}{T}
\simeq
\frac12
\left(
\frac{m_\gamma c_{\rm int}^{2}}
     {\hbar_{\rm int}\omega}
\right)^2.
\]
La chaîne requise est

\[
\text{aller/retour TPK}
\longrightarrow
\text{mémoire}
\longrightarrow
\mu_T
\longrightarrow
m_\gamma.
\]
On n'identifie pas automatiquement \(\varepsilon\), le quotient d'action \(R_{\rm act}\) et \(\mu_T\) : ils appartiennent à des codomaines distincts. La carte nulle du photon ordinaire demeure exacte tant que le gap et son entrelaceur physique ne sont pas construits.

\Needspace{7\baselineskip}
\subsection{Classification opératorielle des secteurs sombres par domaine, lecteur et codomaine physique}
Un secteur sombre ne se définit pas par une masse arbitraire, mais par le noyau de ses couplages. Si \(P_X\) est son projecteur,

\[
P_{\rm em}P_X=0,
\qquad
P_{\rm color}P_X=0,
\]
tandis qu'il peut se produire

\[
P_{\rm grav}P_X\ne0
\]
ou exister un couplage faible. La loi fournit d'abord un opérateur de fibre ; le canal détermine ensuite si son spectre contient un mode stable, métastable ou résonant.

Un photon sombre exige une seconde connexion \(A'\), sa courbure \(F'\) et un entrelaceur cinétique

\[
\mathcal L_{\rm mix}
=-\frac{\varepsilon_{\rm kin}}2F_{\mu\nu}F'^{\mu\nu}.
\]
Si le second secteur acquiert une masse par Stückelberg, Higgs caché ou frontière, son pôle est distinct de celui du photon nul. La grandeur \(\varepsilon_{\rm kin}\) doit être dérivée d'une incidence commune avant la comparaison des observables.

Soient \(P_{\rm weak},P_{\rm em},P_{\rm color}\) les projecteurs sur les canaux de représentation faible, électromagnétique et chromatique, et \(P_X\) le support de la section sombre. Un WIMP exige un recouvrement faible non nul et l'absence de support électromagnétique et chromatique :

\[
P_{\rm weak}P_X\ne0,
\qquad
P_{\rm em}P_X=P_{\rm color}P_X=0.
\]
Le premier produit exprime une incidence de sous-espaces ; un commutateur non nul ne constituerait pas un critère de couplage, puisque deux projecteurs peuvent commuter tout en partageant un support. L'intensité de l'interaction exige ensuite un opérateur de vertex ou un Hamiltonien de couplage.

Sa stabilité exige une valeur propre unitaire du transport ou une parité protectrice. « WIMP » est une condition de représentation et de persistance, et non une masse prédéterminée.

\Needspace{7\baselineskip}
\subsection{Classification de la section neutrinique stérile par fibre invariante, mélange et charge de jauge nulle}
Un neutrino stérile est un vecteur propre du bloc neutre qui satisfait

\[
P_{\rm weak}n_s=0,
\qquad
P_{\rm mass}n_s=n_s.
\]
Il n'est pas synonyme de la quatrième profondeur structurelle du support neutrinique. Son existence exige un noyau faible non nul, une valeur propre de l'opérateur de masse dans ce noyau et un transport de mélange avec les modes actifs. Le critère est opératoriel et peut être réfuté avant l'introduction d'une échelle externe.

\Needspace{7\baselineskip}
\subsection{Classification des sections axionique et pseudoscalaires par caractère topologique et couplage de bord}
Un candidat axionique doit être un mode pseudoscalaire à terme cinétique normalisé,

\[
CP\,a\,CP^{-1}=-a,
\qquad
\mathcal L_{\rm kin}=\frac{Z_a}{2}\,\partial_\mu a\,\partial^\mu a,
\qquad Z_a>0,
\]
couplé à la densité topologique de couleur. Sa masse procède de la courbure du potentiel effectif,

\[
m_a^2=\frac1{Z_a}
\left.
\frac{\partial^2V_{\rm eff}}{\partial a^2}
\right|_{a=a_0}.
\]
Le canal \(C^*\mapsto-C^*\) fournit une orientation impaire, mais ne suffit pas à construire \(V_{\rm eff}\), le terme cinétique ou le projecteur physique. L'espèce n'est définie que lorsque ces applications sont présentes.

\Needspace{7\baselineskip}
\subsection{Réalisations monopolaires et défauts topologiques : classes de charge et stabilité}
Soit \(F\) la courbure d'une connexion \(U(1)\) dans la normalisation intégrale choisie. Un monopôle magnétique se caractérise par la classe

\[
\frac1{2\pi}\int_{S^2}F=n\ne0.
\]
HMT doit réaliser \(n\) comme holonomie nette d'une famille fermée de parcours et démontrer qu'elle survit au raffinement. Une lacune ponctuelle ou une retenue isolée ne constitue pas un monopôle. Si le défaut est dynamique et stable, sa masse est une valeur propre de l'opérateur linéarisé autour de cette classe.

\Needspace{7\baselineskip}
\subsection{Classification des extensions supersymétriques par appariement des représentations et des caractères}
La coexistence de complétions fermionique et bosonique ne démontre pas la supersymétrie. En quatre dimensions, il faut une famille de charges impaires satisfaisant

\[
Q_\alpha:\mathcal F_-\longleftrightarrow\mathcal F_+,
\qquad
\{Q_\alpha,\bar Q_{\dot\beta}\}
=2\sigma^\mu_{\alpha\dot\beta}P_\mu,
\]
qui entrelace les opérateurs de masse :

\[
Q_\alpha\widehat M_-=\widehat M_+Q_\alpha.
\]
L'identité \(Q^2=H\) est réservée à une réduction hermitienne unidimensionnelle explicite ; elle ne remplace pas l'algèbre relativiste précédente.

Ce n'est qu'alors qu'apparaissent des paires spectrales. La brisure doit dériver la séparation des masses sans utiliser la valeur du partenaire recherché. Sélectrons, neutralinos, gluinos ou autres noms demeurent des réalisations possibles, et non des lignes prédites par la seule existence de CAR et de CCR.

\Needspace{7\baselineskip}
\subsection{Classification des sections de dilaton et de radion par modes scalaires de volume et de compactification}
Le dilaton est une fluctuation de la section d'échelle ; le radion, d'une longueur ou d'un volume interne. Leurs actions élémentaires sont

\[
t_0\mapsto e^{\phi_d}t_0,
\qquad
\mathcal V_{\rm int}\mapsto e^{\phi_r}\mathcal V_{\rm int}.
\]
Pour constituer des particules, ils nécessitent un terme cinétique, un potentiel, un projecteur persistant et un pôle. Auparavant, ce sont des degrés collectifs de la carte dimensionnelle. Cette distinction évite de transformer une liberté de normalisation en nouvelle espèce scalaire.

\Needspace{7\baselineskip}
\subsection{Composés et énergie de liaison}
Un état composé \(Y\) de constituants \(X_1,\ldots,X_k\) se construit dans l'ordre

\[
(\Gamma_{X_1},\ldots,\Gamma_{X_k})
\longrightarrow
\mathscr L_{\rm bind}
\longrightarrow
\sigma_Y
\longrightarrow
\eta_Y
\longrightarrow
\widehat M_Y.
\]
L'enregistrement de liaison \(\mathscr L_{\rm bind}\) contient l'ordre des parcours, la frontière commune, la retenue transférée, le canal de couleur et l'holonomie de liaison. Il ne peut être remplacé par \(\sum_i m_{X_i}\). Mésons et baryons exigent des lecteurs distincts parce que leurs nombres de constituants, leurs représentations et leurs symétries d'échange diffèrent.

Les six évaluations composées publiées prouvent l'évaluation du caractère une fois la signature fixée. La clôture état par état exige en outre que l'enregistrement de liaison régénère cette signature. Cette condition s'applique de la même manière aux noyaux, atomes, molécules et candidats tétraquarks ou pentaquarks, avec leurs Hamiltoniens de frontière respectifs.

\Needspace{7\baselineskip}
\subsection{Largeur comme 2e observable}
Les 384 largeurs ne s'obtiennent pas en ajoutant une sixième coordonnée à \(\sigma\). Pour chaque état instable, on construit les canaux \(\mathcal H_{\rm in}\oplus\mathcal H_{\rm out}\), l'opérateur effectif \(H_{\rm eff}(z)\) et le pôle

\[
z_X=M_Xc_{\rm int}^{2}-\frac{i}{2}\Gamma_X.
\]
La partie réelle relève de la loi de masse ; la partie imaginaire, du lecteur de désintégration. La relation \(\tau_X=\hbar_{\rm int}/\Gamma_X\) s'interprète après fixation de la carte temporelle. Une théorie exhaustive de 471 masses et de 384 largeurs contient donc deux lecteurs coordonnés, et non une seule colonne de 855 nombres hétérogènes.

\Needspace{7\baselineskip}
